\Article{7}
Wir beginnen mit der Auswertung des Integrals $\int\mathrm{d}s.\varphi(\mu,\mathrm{d}s)$ über alle Teile des Raumbereichs $s$. Dabei bezeichnet $\mu$ einen bestimmten Punkt außerhalb oder innerhalb des Raumbereichs $s$. Man denke sich eine Kugeloberfläche vom Radius $=1$ um den Mittelpunkt $\mu$ be\SourceLeaf{44}schrieben und in unendlich kleine Elemente zerlegt. Es sei $\mathrm{d}\Pi$ ein solches Element, und eine Gerade von $\mu$ zu einem Punkt dieses Elements schneide die Oberfläche des Raumbereichs $s$ der Reihe nach in den Punkten $p',p'',p'''$ usw. Ihre Anzahl ist gerade oder ungerade, je nachdem, ob $\mu$ außerhalb oder innerhalb des Raumbereichs $s$ liegt. Die Abstände $\mu p',\mu p'',\mu p'''$ usw. bezeichnen wir mit $r',r'',r'''$ usw. Ziehen wir ferner von $\mu$ zu jedem Punkt des Randes des Elements $\mathrm{d}\Pi$ eine Gerade, so entsteht ein pyramidenförmiger Raumbereich. Die Geraden schneiden aus der Oberfläche des Raumbereichs $s$ bei den Punkten $p',p'',p'''$ usw. Elemente heraus, die wir der Reihe nach mit $\mathrm{d}t',\mathrm{d}t'',\mathrm{d}t'''$ usw. bezeichnen. Schließlich sei $q'$ der Winkel zwischen der Geraden $p'\mu$ und der nach außen gerichteten Normalen auf dem Element $\mathrm{d}t'$. Ebenso seien $q'',q'''$ usw. die Neigungen der entsprechenden Normalen bei den Punkten $p'',p'''$ usw. gegen die nach $\mu$ gezogene Gerade. Damit gilt offenbar
\[
\mathrm{d}\Pi=\pm\frac{\mathrm{d}t'.\cos q'}{r'r'}=\mp\frac{\mathrm{d}t''.\cos q''}{r''r''}=\pm\frac{\mathrm{d}t'''.\cos q'''}{r'''r'''}\ \text{etc.}
\]
wobei die oberen oder unteren Vorzeichen gelten, je nachdem, ob $\mu$ außerhalb oder innerhalb des Raumbereichs $s$ liegt.

Ferner ist klar, dass das Integral $\int\mathrm{d}s.\varphi(\mu,\mathrm{d}s)$ über die in jenem pyramidenförmigen Raumbereich enthaltenen Teile des Raumbereichs $s$ durch das Integral $\mathrm{d}\Pi.\int rr\varphi r.\mathrm{d}r$ gegeben ist. Es wird von $r=r'$ bis $r=r''$, dann von $r=r'''$ bis $r=r''''$ usw. genommen, wenn $\mu$ außerhalb des Raumbereichs $s$ liegt, oder von $r=0$ bis $r=r'$, dann von $r=r''$ bis $r=r'''$ usw., wenn $\mu$ innerhalb des Raumbereichs $s$ liegt. Setzen wir also für das unbestimmte Integral
\[
\int rr\varphi r.\mathrm{d}r=-\psi r
\]
mit beliebig gewählter Integrationskonstante, so ergibt sich für das Integral $\int\mathrm{d}s.\varphi(\mu,\mathrm{d}s)$, soweit es über die innerhalb jenes pyramidenförmigen Raumbereichs liegenden Teile des Raumbereichs $s$ erstreckt wird,
\[
\begin{aligned}
&=\mathrm{d}\Pi.(\psi r'-\psi r''+\psi r'''-\ \text{etc.})\\
&=\frac{\mathrm{d}t'.\cos q'.\psi r'}{r'r'}+\frac{\mathrm{d}t''.\cos q''.\psi r''}{r''r''}+\frac{\mathrm{d}t'''.\cos q'''.\psi r'''}{r'''r'''}+\ \text{etc.}
\end{aligned}
\]
wenn $\mu$ außerhalb des Raumbereichs $s$ liegt, dagegen
\[
\begin{aligned}
&=\mathrm{d}\Pi.(\psi0-\psi r'+\psi r''-\psi r'''+\ \text{etc.})\\
&=\mathrm{d}\Pi.\psi0+\frac{\mathrm{d}t'.\cos q'.\psi r'}{r'r'}+\frac{\mathrm{d}t''.\cos q''.\psi r''}{r''r''}+\frac{\mathrm{d}t'''.\cos q'''.\psi r'''}{r'''r'''}+\ \text{etc.}
\end{aligned}
\]
wenn $\mu$ innerhalb des Raumbereichs $s$ liegt.

Wird nun über alle Teile der Kugeloberfläche summiert, so ergibt sich für das vollständige Integral $\int\mathrm{d}s.\varphi(\mu,\mathrm{d}s)$
\SourceLeaf{45}
\[
\begin{aligned}
\text{im ersten Fall}\qquad&=\int\frac{\mathrm{d}t.\cos q.\psi r}{rr}\\
\text{im zweiten Fall}\qquad&=4\pi\psi0+\int\frac{\mathrm{d}t.\cos q.\psi r}{rr}
\end{aligned}
\]
wobei $\mathrm{d}t$ allgemein die Elemente der Oberfläche des Raumbereichs $s$ bezeichnet. Für diese Elemente bezeichnen $q,r$ dieselben Größen wie zuvor die mit Strichen versehenen Buchstaben für bestimmte Elemente. Schließlich bezeichnet $\pi$ den halben Umfang eines Kreises vom Radius $=1$.

Man sieht ferner leicht, dass für einen Punkt $\mu$, der weder außerhalb noch innerhalb des Raumbereichs $s$, sondern auf seiner Oberfläche liegt, die zweite Formel gilt, wenn der Faktor $4\pi$ durch $2\pi$ ersetzt wird und die Oberfläche bei $\mu$ weder eine Spitze noch eine Kante aufweist. Für unseren Zweck brauchen wir diesen Fall aber nicht zu betrachten.

\Article{8}
Nach der Untersuchung des vorigen Artikels führt die Auswertung des Integrals $\int\int\mathrm{d}s.\mathrm{d}S.\varphi(\mathrm{d}s,\mathrm{d}S)$ auf
\[
4\pi\sigma\psi0+\int\int\mathrm{d}t.\mathrm{d}S.\frac{\cos q.\psi(\mathrm{d}t,\mathrm{d}S)}{(\mathrm{d}t,\mathrm{d}S)^2}
\]
wenn wir mit $\sigma$ das Volumen des Raumbereichs bezeichnen, den beide Raumbereiche $s,S$ gemeinsam haben. Der erste Teil $4\pi\sigma\psi0$ fällt damit weg, wenn die Raumbereiche $s,S$ einander ausschließen. Es bleibt ein neues Integral übrig, das weiterhin doppelt erscheint, tatsächlich aber fünffach ist. Um es auf ein vierfaches Integral zurückzuführen, betrachten wir das Integral
\[
\int\mathrm{d}S.\frac{\cos q.\psi(\mu,\mathrm{d}S)}{(\mu,\mathrm{d}S)^2}
\]
über alle Elemente des Raumbereichs $S$. Wieder bezeichnet $\mu$ einen bestimmten Punkt und $q$ den Winkel zwischen zwei von diesem Punkt ausgehenden Geraden, von denen die eine zum Element $\mathrm{d}S$ hin verläuft und die andere fest ist. Wir werden nun auf zwei ganz verschiedene Arten zeigen, wie sich dieses scheinbar einfache, tatsächlich dreifache Integral auf ein wirklich doppeltes Integral zurückführen lässt.

Die durch den Punkt $\mu$ gehende, auf jener festen Geraden senkrechte Ebene bezeichnen wir mit $\Pi$. Man denke sich den Teil dieser Ebene, den die Projektion des Raumbereichs $S$ erreicht, in unendlich kleine Elemente $\mathrm{d}\Pi$ zerlegt. Durch einen Punkt eines solchen Elements $\mathrm{d}\Pi$ ziehe man eine Gerade senkrecht zur Ebene $\Pi$. Sie schneide die Oberfläche des Raumbereichs $S$ beim Fortschreiten in der zur festen Geraden parallelen Richtung der Reihe nach in den Punkten $P',P'',P'''$ usw. Deren Abstände vom Punkt $\mu$ seien der Reihe nach $R',R'',R'''$ usw. Entsprechende Geraden durch alle Punkte des Randes des Ele\SourceLeaf{46}ments $\mathrm{d}\Pi$, rechtwinklig zur Ebene gezogen, bilden einen prismenförmigen Raumbereich. Bei den Punkten $P',P'',P'''$ usw. schneiden sie aus der Oberfläche des Raumbereichs $S$ Elemente heraus, die wir mit $\mathrm{d}T',\mathrm{d}T'',\mathrm{d}T'''$ usw. bezeichnen. Schließlich sei $\chi'$ der Winkel zwischen zwei vom Punkt $P'$ ausgehenden Geraden, von denen die eine die nach außen gerichtete Normale auf dem Element $\mathrm{d}T'$ und die andere zur festen Geraden parallel ist. Die entsprechenden Winkel bei den Punkten $P'',P'''$ usw. seien $\chi'',\chi'''$ usw. Damit gilt offenbar
\[
\mathrm{d}\Pi=-\mathrm{d}T'.\cos\chi'=+\mathrm{d}T''.\cos\chi''=-\mathrm{d}T'''.\cos\chi'''\ \text{etc.}
\]
Man zerlege den prismenförmigen Raumbereich in unendlich kleine Elemente $\mathrm{d}\Pi.\mathrm{d}z$. Dabei bezeichnet $z$ den Abstand eines beliebigen Punkts von der Ebene $\Pi$, positiv auf der Seite gerechnet, nach der die feste Gerade weist. Bezeichnen wir den Abstand desselben Punkts vom Punkt $\mu$ mit $r$, so gilt $z=r\cos q$ und, da $rr-zz$ konstant ist, auch $z\mathrm{d}z=r\mathrm{d}r$ oder $\mathrm{d}\Pi.\mathrm{d}z.\cos q=\mathrm{d}\Pi.\mathrm{d}r$. Daraus folgt, dass unser Integral $\int\mathrm{d}S.\frac{\cos q.\psi(\mu,\mathrm{d}S)}{(\mu,\mathrm{d}S)^2}$ über die im prismenförmigen Raumbereich enthaltenen Teile des Raumbereichs $S$ durch das Integral $\mathrm{d}\Pi.\int\frac{\mathrm{d}r.\psi r}{rr}$ gegeben ist, wenn es von $r=R'$ bis $r=R''$, dann von $r=R'''$ bis $r=R''''$ usw. erstreckt wird. Setzen wir also für das unbestimmte Integral
\[
\int\frac{\mathrm{d}r.\psi r}{rr}=-\vartheta r
\]
mit beliebig gewählter Integrationskonstante, so ergibt sich für unser Integral über die innerhalb des prismenförmigen Raumbereichs liegenden Teile des Raumbereichs $S$
\[
\begin{aligned}
&=\mathrm{d}\Pi.(\vartheta R'-\vartheta R''+\vartheta R'''-\ \text{etc.})\\
&=-\mathrm{d}T'.\cos\chi'.\vartheta R'-\mathrm{d}T''.\cos\chi''.\vartheta R''-\mathrm{d}T'''.\cos\chi'''.\vartheta R'''-\ \text{etc.}
\end{aligned}
\]
Fassen wir diese Summen über die zu sämtlichen Elementen $\mathrm{d}\Pi$ gehörenden Prismen zusammen, so sind offenbar alle Elemente der Oberfläche des Raumbereichs $S$ erfasst, und wir erhalten das vollständige Integral
\[
\int\mathrm{d}S.\frac{\cos q.\psi(\mu,\mathrm{d}S)}{(\mu,\mathrm{d}S)^2}=-\int\mathrm{d}T.\cos\chi.\vartheta R
\]
wobei $\mathrm{d}T$ ein beliebiges Element der Oberfläche des Raumbereichs $S$, $R$ seinen Abstand vom Punkt $\mu$ und $\chi$ den Winkel zwischen der nach außen gerichteten Normalen auf dem Element $\mathrm{d}T$ und einer zur festen Geraden parallelen Geraden bezeichnet.

Auf diese Weise ist das Integral $\int\int\mathrm{d}s.\mathrm{d}S.\varphi(\mathrm{d}s,\mathrm{d}S)$ auf die Form
\[
4\pi\sigma\psi0-\int\int\mathrm{d}t.\mathrm{d}T.\cos\chi.\vartheta(\mathrm{d}t,\mathrm{d}T)
\]
zurückgeführt. \SourceLeaf{47}Dabei bezeichnet $\chi$ offenbar die Neigung der Elemente $\mathrm{d}t,\mathrm{d}T$ gegeneinander, gemessen an der Neigung ihrer bezüglich der Raumbereiche $s,S$ jeweils nach außen gerichteten Normalen gegeneinander. Die Integrationen sind über die vollständigen Oberflächen beider Raumbereiche zu erstrecken.

